\section{为什么需要 Retrieval Augmentation}

课程先把 retrieval 放回 language modeling 的历史，而不是从某个 2023 年产品名开始。这样做能区分三个层次：autoregressive objective 很旧，规模带来新能力，instruction tuning 修复人机接口，而 retrieval 处理的是知识更新、证据来源与可控记忆。

\subsection{语言模型不是近几年才出现}

Age of Language Models slide 用一个简单句子说明 next-token factorization：预测 ``going'' 时，前面的 ``Where are we'' 是 context。下一页展示 1991 年 neural language modelling 工作，提醒读者语言模型与神经网络结合已有数十年历史。

\lecturefigure{slide-02-age-of-language-models.jpg}{语言模型把序列概率分解为逐 token 条件概率}{Stanford 课堂视频 00:02:21}

\begin{importantbox}{Autoregressive factorization}
对 token 序列 $w_{1:T}$，自回归语言模型写成
\[
p(w_{1:T})=\prod_{t=1}^{T}p(w_t\mid w_{<t}).
\]
$w_{<t}$ 是位置 $t$ 之前的上下文，模型只需反复估计下一个 token 的条件分布。LLM 的参数规模和训练数据巨大，但这个概率分解本身并不新。
\end{importantbox}

\lecturefigure{slide-03-elman-1991-neural-lm.jpg}{课堂展示的 1991 年 neural language modelling 论文}{Stanford 课堂视频 00:02:57}

\begin{knowledgebox}{读图：历史页服务于概念校准}
这张扫描页不是要求记住旧网络结构，而是提醒：distributed lexicon、neural sequence processing 与概率语言建模早已存在。今天的突破来自规模、优化、硬件、数据和 interface 的共同演化，不能把整段历史缩成某家公司“发明语言模型”。
\end{knowledgebox}

\teachervoice{Kiela 开玩笑说，如果认为 OpenAI 发明了语言模型，他会生气。这个幽默在教学上承担了校准作用：autoregressive idea 很旧，LLM 的新颖性主要来自 scale 与现代训练栈。}

\subsection{Next-token prediction 与“坏掉的用户界面”}

Next Word Prediction slide 把系统画成 input--generator--output，并指出 plain sequence-to-token 的问题之一是 broken user interface。用户过去需要构造奇怪 prompt 才能让模型执行任务；ChatGPT 式 instruction tuning 与 alignment 让人能够直接表达意图。

\lecturefigure{slide-04-next-word-prediction.jpg}{Plain next-word prediction 的简单结构与界面问题}{Stanford 课堂视频 00:04:30}

\begin{knowledgebox}{读图：模型 objective 与用户 objective 不同}
预训练只要求预测语料中的下一 token，用户却希望模型遵循命令、拒绝危险请求、保持格式并完成任务。Prompting、instruction fine-tuning 与 RLHF 是 interface alignment：它们让模型更容易被调用，但没有自动解决知识陈旧、证据缺失或事实错误。
\end{knowledgebox}

\teachervoice{讲者把 ChatGPT 的关键贡献概括为“fix the user interface”。人更擅长直接说想要什么，而不是猜测哪个奇怪 prompt 能触发行为；instruction data 在普通网页语料中又非常稀少。}

\subsection{把输出 elicitation 修好之后，仍然有五类缺口}

Eliciting Outputs slide 把 prompt、instruction tuning 与 alignment 放在 generator 之前，右侧列出 hallucination、attribution、staleness、revisions 与 customization。这五项不是同一个错误率，而是分别对应事实、来源、时间、可更新性与用户/组织边界。

\lecturefigure{slide-05-eliciting-outputs.jpg}{输出 elicitation 改善后仍然存在的模型缺口}{Stanford 课堂视频 00:06:09}

\begin{knowledgebox}{术语消化：五类问题不能混成“准确率”}
\begin{tabularx}{\textwidth}{L{0.18\textwidth}>{\raggedright\arraybackslash}X >{\raggedright\arraybackslash}X}
\toprule
问题 & 表现 & Retrieval 可能提供什么 \\
\midrule
Hallucination & 输出与可接受事实源冲突 & 提供显式证据，但生成器仍可能忽略 \\
Attribution & 无法指出来源 & 返回文档 ID、段落和时间戳 \\
Staleness & 参数知识过时 & 替换或刷新 index，而不重训全部模型 \\
Revision & 错误事实难以精确更正 & 修改文档或访问策略 \\
Customization & 不同组织需要不同知识边界 & 切换私有/领域 index 与权限 \\
\bottomrule
\end{tabularx}
\end{knowledgebox}

\begin{warningbox}{检索到证据不等于 grounded generation}
Retriever 可能找到正确文档，generator 仍可能引用错段落、混入参数记忆或完全忽略证据。因此至少要分别评估 retrieval quality、evidence utilization 与 final answer correctness。
\end{warningbox}

\teachervoice{讲者强调，尤其在 strict accuracy requirements 的 enterprise 场景，模型高置信地“make up stuff”、不给 attribution、不能快速修订和定制，都会阻止真实部署。}

\subsection{Contextualization：外部记忆进入生成器}

Contextualization architecture 增加 documents、document encoder、retriever、query encoder 与 context。用户输入不再只进入 generator，也形成 query；retriever 从外部文档集合中找证据，再把结果与 prompt 一起交给模型。

\lecturefigure{slide-06-contextualization-architecture.jpg}{Documents--retriever--context--generator 的基础 RAG 架构}{Stanford 课堂视频 00:06:48}

\begin{knowledgebox}{第一使用语表：Retriever、Index、Embedding、Context}
\begin{tabularx}{\textwidth}{L{0.17\textwidth}>{\raggedright\arraybackslash}X}
\toprule
术语 & 含义 \\
\midrule
Document chunk & 从文档切出的可检索单位，粒度决定召回与上下文成本 \\
Embedding & 把 query 或 chunk 映射为向量，以便计算相似度 \\
Index & 支持查找的数据结构，可为倒排索引、向量索引或混合索引 \\
Retriever & 根据 query 从 index 返回候选 chunk 与分数 \\
Context & 最终提供给 generator 的证据文本，不等同于整个 index \\
\bottomrule
\end{tabularx}
\end{knowledgebox}

\lecturefigure{slide-07-two-paradigms.jpg}{Closed-book/open-book 与 parametric/semi-parametric 两组范式}{Stanford 课堂视频 00:07:48}

\begin{importantbox}{Parametric 与 non-parametric memory}
Parametric memory 是模型权重中通过训练形成的统计知识，更新通常需要继续训练；non-parametric memory 是可直接读写或替换的文档、数据库、索引与缓存。RAG 常被称为 semi-parametric，因为生成器有参数，而外部 memory 可在测试时变化。
\end{importantbox}

\subsection{为什么 external index 能缓解部分问题}

Why does this solve the issues slide 给出两条直接收益：替换 index 可以 customization、避免 staleness、修订错误；grounding 可以减少 hallucination，并提供 citation 与 attribution。注意用词是“less”而非“zero”。

\lecturefigure{slide-08-why-retrieval-solves-issues.jpg}{外部 index 对定制、更新与 grounding 的帮助}{Stanford 课堂视频 00:08:57}

\begin{knowledgebox}{读图：可更新性来自状态位置变化}
若事实只存在于权重中，修正需要训练与重新部署；若事实位于 index，更新路径变成 document update $\rightarrow$ re-index $\rightarrow$ 下一次 retrieval。代价是系统必须管理文档版本、权限、chunking 和检索一致性。
\end{knowledgebox}

\begin{warningbox}{外部知识库不是自动真相库}
Index 可能包含错误、冲突、过期或恶意内容。RAG 把“模型是否记错”部分转化成“系统信任什么数据源、如何排序与引用”；它提高可控性，却把数据治理和 source selection 变成显式责任。
\end{warningbox}

\teachervoice{课堂提醒，index 可换、事实可修订、来源可引用，但 generator 仍可能不理会 context。检索增强真正要求 retriever 与 generator 之间建立有效接口。}

\subsection{一张架构图会产生许多研究问题}

Many Questions slide 围绕基础图追问：系统如何学习、扩展与预处理，怎样编码和检索文档，又该如何设计 prompt 并验证答案。它把 RAG 从单一 trick 展开为完整的系统设计空间，也为后面按 train/test、representation 和 fusion 分类做准备。

\lecturefigure{slide-09-many-questions.jpg}{RAG 基础架构周围的学习、扩展、检索与验证问题}{Stanford 课堂视频 00:09:57}

\begin{importantbox}{RAG 的六个设计轴}
\begin{enumerate}
  \item \textbf{Corpus}：哪些文档可被访问，谁维护版本与权限；
  \item \textbf{Chunking}：如何切分、重叠、保留层级；
  \item \textbf{Representation}：sparse、dense、late interaction 或 hybrid；
  \item \textbf{Retrieval policy}：何时检索、取多少、是否多跳；
  \item \textbf{Fusion}：拼 prompt、cross-attention、decoder fusion 或 token-level memory；
  \item \textbf{Learning/evaluation}：哪些组件有梯度，怎样区分检索与生成错误。
\end{enumerate}
\end{importantbox}

\subsection{本章小结}

Retrieval augmentation 主要补充可更新、可归因和可定制的外部记忆。它没有改变 next-token objective，也不会自动保证生成忠实；因此需要把 corpus、index、retriever、fusion、generator 和 evaluation 当作一套系统来设计。

